% =========================================================
% Protocolo de la prueba integrada con el código congelado
% =========================================================

\subsection{Protocolo de la prueba integrada}
\label{sec:protocolo_prueba}

La prueba integrada contrasta, sobre una misma muestra y un mismo
montaje, la captura realizada con el método habitual del laboratorio
---el osciloscopio operado desde su panel frontal, sin la
aplicación--- con la que entrega el sistema automatizado en su versión
congelada (rama \texttt{main}, commit \texttt{c67ae06}). La muestra es
agua destilada. Al igual que en el resto de la \cref{sec:prueba_integrada},
no se pretende caracterizar el líquido: el objeto de la comparación es
la herramienta, y la pregunta es si el sistema reproduce lo que el
método habitual obtiene y si, además, conserva la información necesaria
para atribuir cada dato a sus condiciones.

\subsubsection*{Condiciones que deben ser idénticas}

La comparación solo es válida si la única diferencia entre ambos
conjuntos de capturas es el método de adquisición. Se fijan, por tanto,
antes de la primera captura y sin modificarse hasta la última:

\begin{itemize}[leftmargin=*, itemsep=2pt]
  \item \textbf{Muestra y celda.} El mismo volumen de agua destilada en
    la misma celda, sin moverla ni rellenarla entre métodos.
  \item \textbf{Geometría.} La separación hidrófono--haz se mide una
    sola vez. Se captura, junto con el volumen, en la sección de
    geometría de la aplicación antes de crear la sesión, porque al
    crearse los campos se bloquean y ya no pueden añadirse.
  \item \textbf{Láser.} Los mismos valores de \textit{Output level},
    \textit{EO delay}, \textit{Burst mode} y temperatura de
    enfriamiento en ambos métodos, anotados al inicio. Al terminar cada
    secuencia automática la aplicación envía el láser a modo seguro
    (\cref{tab:modo_seguro}), que fija \textit{E OFF} y un
    \textit{EO delay} de \SI{3800}{\micro\second}; si la prueba
    continúa después, los parámetros acordados deben restablecerse y
    comprobarse antes de la siguiente captura.
  \item \textbf{Osciloscopio.} El mismo canal para el hidrófono, las
    mismas escalas vertical y horizontal, la misma posición horizontal
    y el mismo modo y nivel de disparo. La aplicación no modifica estos
    ajustes; los lee y los registra en cada fila. La posición horizontal
    fija el origen del eje de tiempos (\cref{sec:origen_tiempos}), de
    modo que cualquier cambio invalida la comparación del instante de
    llegada.
  \item \textbf{Promediado.} El mismo modo de adquisición y el mismo
    número de promedios en ambos métodos. La secuencia automática por
    tiempo impone el modo \textit{Average} con 100 promedios durante su
    ejecución y restaura el modo previo al terminar; si se emplea, el
    método habitual debe usar ese mismo promediado. Si se prefiere otro,
    las capturas del sistema deben tomarse en modo manual, que respeta
    la configuración del panel.
  \item \textbf{Número de capturas.} El mismo número en ambos métodos,
    fijado antes de comenzar.
  \item \textbf{Temperatura.} Se anota la temperatura de la muestra al
    inicio y al final de cada método. En el método habitual es la única
    referencia térmica; en el sistema, cada fila registra además la
    lectura de los cuatro sensores.
\end{itemize}

\subsubsection*{Secuencia de la prueba}

\begin{enumerate}[leftmargin=*, itemsep=2pt]
  \item Se prepara el montaje, se ajustan escalas y disparo en el
    osciloscopio y se anotan en la hoja de sesión los parámetros del
    láser, la geometría, la hora de la computadora y la del
    osciloscopio, cuyos relojes pueden diferir.
  \item Se adquiere el conjunto del método habitual y se conservan los
    archivos que guarda el osciloscopio.
  \item Sin tocar el montaje ni el osciloscopio, se abre la aplicación,
    se verifica que los tres indicadores de conexión estén en verde y se
    capturan los datos de geometría.
  \item Se realiza y guarda una captura manual, que crea la sesión y
    permite comprobar la señal; a continuación se ejecuta la secuencia
    por tiempo con el número de capturas acordado y un intervalo de al
    menos \SI{15}{\second}. La aplicación rechaza el intervalo si no
    alcanza para completar el promedio y transferir la forma de onda.
  \item Si el tiempo disponible lo permite, se repite el método habitual
    al final. Un conjunto habitual antes y otro después del automatizado
    permite distinguir una deriva del montaje durante la prueba de una
    diferencia atribuible al método.
  \item Se copia la carpeta completa de la sesión, junto con los
    archivos del método habitual, la hoja de sesión y el registro de
    incidencias.
\end{enumerate}

\subsubsection*{Criterios de comparación}

Los criterios se limitan a magnitudes que el sistema ya mide o
registra, y se calculan con el mismo estimador sobre ambos conjuntos:
el que implementa la herramienta \texttt{verificar\_sesion.py} del
directorio \texttt{herramientas} del repositorio. La traza
se reconstruye con las \cref{eq:conv_voltaje,eq:conv_tiempo} y con los
parámetros de preámbulo propios de cada archivo, se le resta la línea
base estimada al inicio del registro y sobre ella se evalúan las
magnitudes siguientes.

\begin{enumerate}[leftmargin=*, itemsep=2pt]
  \item \textbf{Amplitud pico a pico.} Para cada método se obtienen la
    media del valor pico a pico sobre la traza completa y su
    coeficiente de variación, como en la \cref{sec:repetibilidad}. La
    diferencia entre las medias de ambos métodos se expresa en relación
    con la mayor de las dos dispersiones internas: una diferencia del
    orden de esa dispersión indica que el sistema reproduce el método
    habitual; una diferencia que la excede debe explicarse antes de
    atribuirla al método.
  \item \textbf{Instante de llegada del máximo.} Para cada captura se
    obtiene el instante del máximo de la amplitud absoluta en el eje
    reconstruido, y se comparan su media y su dispersión entre métodos.
    Con escalas y posición horizontal idénticas, ambos conjuntos deben
    situar el máximo en el mismo instante dentro de su dispersión.
  \item \textbf{Banderas de error.} Se cuenta cuántas filas del sistema
    tienen \texttt{error\_flag = 1} y se lee el \texttt{error\_desc} de
    cada una. Las filas marcadas se excluyen de los dos criterios
    anteriores, pero se reportan con su causa: que el sistema marque
    una captura dudosa en lugar de omitirla es parte de lo que se
    evalúa.
  \item \textbf{Metadatos.} Se comprueba que cada fila del sistema
    contenga poblados los campos de temperatura y sensores, preámbulo,
    escalas, modo de adquisición, promedios solicitados y
    efectivamente promediados, y parámetros del láser
    (\cref{sec:formato_datos}), y que \texttt{metadatos\_sesion.txt}
    registre la identidad del osciloscopio y la geometría. Se contrasta
    con lo que, para el método habitual, solo consta en la hoja de
    sesión.
  \item \textbf{Comprobaciones de sanidad.} Sobre ambos conjuntos se
    verifica que ninguna captura sature el convertidor, que no existan
    pares de registros idénticos bit a bit y que todas las capturas de
    un mismo conjunto compartan la configuración de escala; son las
    alertas que emite la misma herramienta.
\end{enumerate}

La prueba no fija de antemano un umbral de aceptación para la diferencia
entre métodos: la referencia es la dispersión interna que la propia
prueba mida. Lo que sí se fija de antemano es que ambos conjuntos se
procesen con el mismo estimador y que se reporten todas las capturas,
incluidas las marcadas.
